\section{The HBZ/rANS dynamical system and fixed-input failure exclusion}
\label{sec:rans-frontier}

\subsection{Why success postconditions and termination are different claims}

The inverse function for the signature prefix, the HBZ coefficient--symbol leaf,
the edge mapping canonical symbol streams through the actual encoder to pure byte
traces, the edge mapping exact traces through the actual decoder back to canonical
symbol streams, core composition across an actual copy, and the full-HBZ wrapper
boundary of production SignaturePack/Unpack have each been closed.
Week~15 excluded failure of the actual encoder when it returns on the all-6
stream produced by canonical all-zero HBZ, and lifted this to a production zero
round trip \cite{week15-report,rans-actual-success-witness}. Thus
\claimid{OBL-RANS-ACTUAL-SUCCESS-WITNESS} is \statusproved\ as fixed-input Hoare
partial correctness. It is not a termination/losslessness theorem guaranteeing
a return or the existence of an execution. After this result, the project moved
through the frozen Week~16 KeyGen/Sign/Verify \FrozenPaperGate\ to the
\CurrentGate\ of the post-freeze KeyGen semantics audit.

The integer arithmetic of rANS itself is not complicated. The difficulty lies
in implementation direction: the encoder reads symbols backwards and decrements
the byte-buffer cursor, while the decoder reads the copied buffer forwards.
The key is an invariant connecting the pure inverse-function theorem to the
states of the two actual nested loops. Week~11 completed the reverse-encoder
connection, Week~12 the forward-decoder connection, and Week~13 the actual copy
and premise transport between the two theorems. Week~14 composed this core
theorem with the prepare/apply, frame and size semantics of production
\identifier{_encode_hb_z1_full}/\identifier{_decode_hb_z1_full}.
Week~15 combines the fact that actual prepare produces an all-6 stream, a
1020-byte normalization budget, and a necessary condition for actual failure
to force the nonzero-size branch in returning fixed-input executions.
The remaining HBZ gaps are termination and existence of this execution itself,
global success on arbitrary canonical inputs, malformed rejection, and
probabilistic losslessness.

\subsection{The 13-symbol partition}

Mapping a canonical HBZ coefficient \(h_i\in[-6,7)\) to
\(s_i=h_i+6\) gives the alphabet
\[
  \Sigma=\{0,1,\ldots,12\}
\]
Let the rANS scale be \(M=1024\), and let each symbol's cumulative starting
point and frequency be
\[
\begin{aligned}
(c_s)_{s=0}^{12}
  &=(0,1,2,3,8,66,312,710,957,1016,1021,1022,1023),\\
(f_s)_{s=0}^{12}
  &=(1,1,1,5,58,246,398,247,59,5,1,1,1)
\end{aligned}
\]
The certificate establishes
\[
  f_s>0,\qquad \sum_{s=0}^{12}f_s=M,\qquad
  [c_s,c_s+f_s)\text{ partition }[0,M)\text{}
\]

This finite partition uniquely recovers the symbol from the residue slot
\(y\bmod M\). The fact that the actual encoder/decoder's packed tables represent
exactly these \(c_s,f_s\) values and the symbol lookup is also certified against
the generated literal arrays.

\subsection{The pure step is an inverse of Euclidean division}

For a fixed symbol \(s\), define
\begin{align*}
  E_s(x)
    &=\left\lfloor\frac{x}{f_s}\right\rfloor M
      +c_s+(x\bmod f_s),\\
  D_s(y)
    &=\left\lfloor\frac{y}{M}\right\rfloor f_s
      +(y\bmod M)-c_s
\end{align*}
Euclidean division
\(x=f_s\lfloor x/f_s\rfloor+(x\bmod f_s)\) and
\(0\le x\bmod f_s<f_s\) immediately suggest
\[
  D_s(E_s(x))=x
\]
The EasyCrypt theorems
\theoremname{pure_rans_step_quotient},
\theoremname{pure_rans_step_slot}, and
\theoremname{pure_rans_step_inverse} prove these calculations separately.

The actual encoder uses a fast formula with table reciprocals, shifts and biases
instead of division. A separate certificate shows that under
\[
  1\le x<x_{\max}(s),\qquad
  x_{\max}(s)=2^{21}f_s
\]
the fast quotient equals the Euclidean quotient, so \(D_s\) also inverts the
fast step. This is not merely a computational check that ``the table values look
right''; it is a theorem connecting the integer meanings of the actual word fields.

\begin{verified}[Pure and fast steps]
For canonical symbols \(0\le s<13\) and the specified state ranges,
\[
  D_s(E_s(x))=x,\qquad
  D_s(E^{\mathrm{fast}}_s(x))=x
\]
compile successfully. There is also a concrete witness to the premises at
\(s=0,x=1\), so the theorem is not vacuous.
\end{verified}

\subsection{Normalization bytes are coordinates of the discarded quotient}

When the state exceeds \(x_{\max}(s)\), divide it by the byte radix
\(B=256\) 0, 1 or 2 times to reduce it before applying the fast step.
The pure definitions are
\[
  \ell_s(x)\in\{0,1,2\},\qquad
  r_s(x)=\left\lfloor\frac{x}{B^{\ell_s(x)}}\right\rfloor
\]
and the low bytes discarded by division are retained in a list
\(\beta_s(x)\) in decoder read order.

The key readback proposition is
\[
  \operatorname{ReadBytes}(r_s(x),\beta_s(x))=x.
\]
For one byte, use \(x=B\lfloor x/B\rfloor+(x\bmod B)\); for two bytes, apply
this identity twice. The actual encoder writes the low byte first while
decrementing the cursor, so the final forward segment orders the two bytes as
\([\mathsf{high},\mathsf{low}]\). Missing this reversal can leave a pure theorem
correct while making it inconsistent with the actual decoder.

\sourcepath{Mode2RansNormalization.ec} provides leaf theorems showing that the
actual \(W32\) shift/truncate/or operations and \(W64\) cursor decrement have
the same meanings as these integer operations.

\subsection{A pure trace for the complete symbol list}

Let the symbol list be \(s_0,\ldots,s_{n-1}\), with \(n=1024\), and let the
final state be \(x_n=L=2^{23}\). Define recursively, so that the encoder proceeds
backwards,
\[
\begin{aligned}
  \ell_j&=\ell_{s_j}(x_{j+1}),\\
  r_j&=\left\lfloor x_{j+1}/256^{\ell_j}\right\rfloor,\\
  x_j&=E^{\mathrm{fast}}_{s_j}(r_j),\\
  \beta_j&=\beta_{s_j}(x_{j+1})
\end{aligned}
\]
The decoder-facing byte string is
\[
  \mathcal B=\operatorname{LE32}(x_0)
    \Vert\beta_0\Vert\cdots\Vert\beta_{n-1}
\]
Define the cursor cuts by
\[
  c_0=4,\qquad c_{j+1}=c_j+|\beta_j|
\]

The project names \(xs/cuts/bytes\) refer respectively to
\((x_0,\ldots,x_n)\), \((c_0,\ldots,c_n)\), and \(\mathcal B\).
These three objects are defined in \sourcepath{Mode2RansByteStack.ec}.
Its \identifier{valid_rans_trace} constructs them directly from the canonical
symbol list rather than assuming arbitrary witnesses.

\begin{verified}[Inductive inversion of the trace]
For a canonical list,
\[
  2^{23}\le x_j<2^{31}
\]
holds for every \(j\). Moreover, the residue slot of \(x_j\) selects \(s_j\),
\(D_{s_j}(x_j)=r_j\), and reading \(\beta_j\) recovers \(x_{j+1}\).
\end{verified}

The mathematics required for the pure decoder's list-level inductive proof is
therefore available. The refinement showing that actual encoder states represent
these trace coordinates was closed in Week~11; the refinement showing that actual
decoder states consume the same coordinates forwards was closed in Week~12.

\subsection{The six encoder targets audited in Week 10}

Week~10 decomposed the required outer-loop invariant into distinct transport
problems and added \WeekTenTargetCount\ files to the authored target manifest.
All six files compile freshly with \identifier{-no-eco}, but completion of each
local proposition must be distinguished from completion of the parent refinement.

\begin{itemize}
  \item \sourcepath{Mode2RansArrayListBridge.ec}: transports the actual byte
        array's symbol stream to a pure list and suffix, and expresses segment
        agreement and the prefix frame. The array--list bridge and one-symbol
        trace recurrence are \statusproved.
  \item \sourcepath{Mode2RansEncoderWordStep.ec}: connects actual word
        expressions using packed reciprocals, shifts and complements to the
        integer meaning of the pure fast step. Word-level equalities for the
        literal mode~2 table are \statusproved.
  \item \sourcepath{Mode2RansEncoderInnerProgress.ec}: represents the pure
        0/1/2-step progress of the normalization inner loop through the state,
        cursor and byte segment already written. This pure lemma is \statusproved.
  \item \sourcepath{Mode2RansEncoderSerialization.ec}: separates the four
        little-endian bytes of the final \(W32\) state, the exterior frame and
        successful size arithmetic. This leaf proposition is \statusproved.
  \item \sourcepath{Mode2RansEncoderTrace.ec}: supplies lemmas for finite
        phases, byte prepending, segments, prefix frames and absence of wrap
        used by the actual inner loop.
  \item \sourcepath{Mode2RansEncoderActualInner.ec}: installs these invariants
        in the actual generated \identifier{_rans_encode} inner loop and derives
        successful offset and size ranges through a separate actual Hoare theorem.
\end{itemize}

\begin{caution}[Local \statusproved\ and the then-\statuspartial\ parent in Week 10]
At the Week~10 baseline, these files certified only lemmas and local actual
partial correctness. Equality of the complete actual suffix with the pure
\identifier{trace_bytes} was absent, so
\claimid{OBL-RANS-ENCODE-REFINEMENT} was \statuspartial. The next section closes
exactly this parent edge in Week~11. These Week~10 targets did not certify success
branch reachability or termination. Fixed all-6 failure exclusion is a separate
Week~15 theorem, and termination remains open.
\end{caution}

\subsection{The actual encoder invariant closed in Week 11}

Let \(i\) be the actual reverse encoder's outer-loop index, \(x\) its current
word state, \(off\) the buffer cursor decreasing from the end, and \(enc\) the
buffer. With the pure trace
\[
  T_i=\operatorname{EncodeTrace}
       (\operatorname{SymbolSuffix}(symbols_0,i))=(x_i,b_i)
\]
the success branch's live invariant consists of the following five conditions
\cite{rans-encoder-invariant}.
\[
\begin{gathered}
  0\le i\le1024,\qquad x=x_i,\qquad off+|b_i|=1024,\\
  \operatorname{SegmentMatches}(enc,off,b_i),\qquad
  \operatorname{PrefixFrame}(enc_0,enc,off),\qquad 4\le off.
\end{gathered}
\]

Week~11 added \WeekElevenTargetCount\ targets. The tail-aware inner invariant
preserves the initial array \(enc_0\) and
\(\operatorname{EncodeTrace}(tail).bytes\) across actual byte writes and recovers
the exact \(0..2\) normalization length when the actual guard exits. Generated
word-step lemmas identify the word update after actual table loads and protect
erasure with the pure fast step; the final four generated stores prepend the
serialized state to the accumulated suffix.

Consequently, the actual procedure's success postcondition compiled as
\[
\begin{gathered}
  \operatorname{SegmentMatches}
  \bigl(enc',off',\operatorname{TraceBytes}
  (\operatorname{SymbolList}(symbols_0))\bigr),\\
  off'+\left|\operatorname{TraceBytes}
  (\operatorname{SymbolList}(symbols_0))\right|=1024,\\
  \operatorname{PrefixFrame}(enc_0,enc',off')
\end{gathered}
\]

\begin{verified}[\claimid{OBL-RANS-ENCODE-REFINEMENT}]
\theoremname{actual_rans_encode_trace_closure} shows that if the actual generated
\identifier{_rans_encode} returns, it either fails or, on success, satisfies all
three relations above, the offset/size ranges and \(bad=0\).
\theoremname{actual_rans_encode_trace_refinement} explicitly exposes equality of
the entire returned suffix at the parent refinement surface. Success is not a
premise; it is one branch of the conclusion determined by the actual return value.
\end{verified}

\begin{caution}[Limits of partial correctness]
This theorem alone does not establish actual encoder termination, success
reachability or an all-6 symbol success witness. In Week~11,
\claimid{OBL-RANS-ACTUAL-SUCCESS-WITNESS} was \statuspartial; the separate
Week~15 theorem excludes only failure in returning fixed all-6 executions.
\end{caution}

\subsection{The actual decoder invariant closed in Week 12}

For actual forward decoder index \(i\), state \(x\) and read cursor
\(off_d\), Week~12 installed the following invariant directly in the generated
outer/inner loops \cite{week12-report,rans-decoder-invariant}.

\begin{enumerate}
  \item The symbol prefix already output equals \(s_0,\ldots,s_{i-1}\);
  \item the current state is \(x_i\) and the current cursor is \(c_i\);
  \item table lookup selects \(s_i\) and reads the normalization segment
        exactly from \([c_i,c_{i+1})\);
  \item the frame of the unwritten tail of the output array is preserved;
  \item at exit, \(x=2^{23}\), \(off_d=size\), and \(bad=0\).
\end{enumerate}

The historical \theoremname{actual_rans_decode_mode2_control} fixed only the
actual tables/state fields and \(bad\in\{0,1\}\). The new input predicate
\identifier{actual_mode2_decoder_trace_input} combines the expected symbol
array, exact \identifier{trace_bytes}, pointwise relations for every
normalization read, state fields \((1024,encoded\_size,13)\), and actual
generated tables.

\begin{verified}[\claimid{OBL-RANS-DECODE-REFINEMENT}]
\theoremname{actual_rans_decode_trace_refinement} directly opens the actual
generated \identifier{_rans_decode} and, under this exact-trace input, recovers
all 1024 symbols and concludes \(off=encoded\_size\), \(bad=0\), and the decoded
tail frame. \theoremname{decoder_outer_trace_exit_components} internally derives
\(x=2^{23}\), also satisfying the actual final-state rejection check.
\end{verified}

This is Hoare partial correctness and does not establish decoder termination.
The exact-trace input is non-circular; Week~13 added a separate composition
theorem showing that the actual harness produces it.

\subsection{Actual harness core composition closed in Week 13}

The harness in \sourcepath{Mode2RansActualInverse.ec} is not an abstract assembly
diagram: it calls three actual generated procedures in the following order.
The theorem in that file closes only control of these calls; the Week~13
semantic closure is separate, in \sourcepath{Mode2RansCoreActualInverse.ec}.
\[
\mathsf{RansEncode}
\longrightarrow
\mathsf{CopyEncodedSuffix}
\longrightarrow
\mathsf{RansDecode}.
\]
The latter two calls run only when the encoder returns \(bad=0\).

The existing harness control conclusion is
\[
\begin{aligned}
  &\mathit{encoderBad}\in\{0,1\},\\
  &\mathit{decoderRan}\Longleftrightarrow\mathit{encoderBad}=0,\\
  &\mathit{decoderRan}\Longrightarrow
    (count,size,m)=(1024,\mathit{encodedSize},13)
\end{aligned}
\]
Week~13's \sourcepath{Mode2RansCoreCompositionBridge.ec} composes the encoder's
\identifier{segment_matches} and the copy's \identifier{slice_eq} into an exact
trace in the copied buffer, and transports it to the decoder's pointwise-read
and configured-state input predicate. In particular, it derives absence of wrap
from actual success bounds instead of silently replacing modular (W64)
subtraction with integer subtraction.

\begin{verified}[\claimid{OBL-RANS-CORE-INVERSE}]
\theoremname{actual_rans_encode_copy_decode_inverse} does not assume actual
encoder success. On failure, the actual copy and decoder are not run, and the
initial decoded array/state is preserved. On success, the following semantic
conclusions appear in the postcondition of the same harness theorem:
\[
  \mathit{decoderBad}=0,\quad
  off_{\mathrm{final}}=size,\quad
  decoded[0..1024)=symbols[0..1024).
\]
The decoder subproof also derives the internal exit state \(x=2^{23}\), but the
harness ABI does not expose it. The decoded tail \([1024,2048)\) is preserved
by a frame. This result is success-conditioned partial correctness.
\end{verified}

At the end of Week~13, no EasyCrypt theorem forced the actual success branch.
A concrete pure-trace witness, a canonical all-zero precondition witness, or a
size-4 decoder-state witness did not replace the postcondition that the actual
encoder returns \(bad=0\) on the concrete all-6 input. Week~15 closes exactly
this gap by contradicting the budget with the actual failure cause.

\subsection{Correctly reading the full HBZ wrapper closed in Week 14}

Because the actual encoder can express failure, canonical \(h\) alone must not
be used to assert an unconditional round trip. The parent proposition proved
in Week~14 must expose the failure branch, as in
\[
  size=0\quad\lor\quad
  \bigl(size\ne0\land bad=0\land
        decoded[0..1024)=h[0..1024)\bigr)
\]
Nonzero size bounds and buffer frames are also required.

\begin{verified}[\claimid{OBL-SIG-HBZ-ENCODE-DECODE}]
\theoremname{actual_hbz_full_encode_decode_inverse_mode2} proves this
failure/success disjunction for the focused actual full-HBZ harness.
\theoremname{signature_pack_unpack_hbz_full_actual_exact} shows that the
production SignaturePack/Unpack harness agrees exactly with the focused
harness in its return tuple and \identifier{decoder_ran} flag, and
\theoremname{signature_pack_unpack_hbz_full_inverse_mode2} lifts the same
contract to the production boundary. The success branch includes
\(4\le size\le1024\), coefficient-prefix recovery, coefficient/suffix frames,
and a prepared-symbol trace witness.
\end{verified}

\subsection{Week 15's fixed all-6 budget and actual failure exclusion}

Opening byte projections to prove that all actual 32-bit loads from
\identifier{zero_hbz} are 0 shows that actual prepare returns \(bad=0\) and a
prefix of 1024 symbol~6 entries. In the literal table, \(f_6=398\) and
\(x_{\max}(6)=2^{21}\cdot398\);
\theoremname{symbol6_normalization_len_le1} shows that each symbol~6 step with
\(2^{23}\le x<2^{31}\) emits at most one normalization byte.

The first four states computed from the initial state are
\[
  8388608\to21582496\to55528910\to142868116\to367580518
\]
and all four steps emit 0 bytes. Thus, even allowing one byte for each of the
remaining 1020 steps gives normalization length at most 1020; adding the final
four state bytes puts the total trace in \([4,1024]\). This is a conservative
capacity argument, not a claim that the observed exact 174/178-byte execution
has been proved as a theorem.

\begin{verified}[\claimid{OBL-RANS-ACTUAL-SUCCESS-WITNESS}]
\theoremname{actual_rans_encode_failure_trace_cause} states that a nonzero
\(bad\) from the actual encoder implies trace length \(>1020\). Contradicting
this with the all-6 upper bound lets
\theoremname{actual_rans_encode_all_six_success} establish \(bad=0\) and the
size/frame/exact-trace postconditions for every returning actual core execution.
\theoremname{actual_encode_hb_z1_full_zero_success},
\theoremname{signature_pack_hbz_zero_success_mode2}, and
\theoremname{signature_pack_unpack_hbz_zero_success_mode2} lift this conclusion
to the focused full wrapper including actual prepare and to the production
SignaturePack/Unpack zero round trip.
\end{verified}

\begin{caution}[\identifier{GO-WITNESS} is not \identifier{GO-WITNESS-LOSSLESS}]
These are Hoare partial-correctness theorems excluding failure in every
\emph{returning} fixed-input execution. Fixed-input termination, execution
existence (non-vacuity), losslessness, probability-one success and
\identifier{phoare} have not been proved. The \(h\) codec became
\statusdeferred\ because work moved to Week~16 MINCORE KeyGen, not because
the success witness was incomplete.
\end{caution}

\begin{center}
\begin{tabularx}{\textwidth}{Y L{31mm}}
\toprule
rANS layer & Status \\
\midrule
13-symbol intervals/frequencies and actual table certificate & \statusproved \\
pure quotient/slot/fast-step inverse & \statusproved \\
0/1/2-byte normalization and W32/W64 leaf semantics & \statusproved \\
canonical \(xs/cuts/bytes\) trace and state bounds & \statusproved \\
actual array/list suffix and word-step bridge & \statusproved \\
actual four-byte final serialization composition & \statusproved \\
actual encoder successful offset/size ranges & \statusproved\ (partial correctness) \\
actual inner normalization byte/frame/no-underflow &
\statusproved\ (tail-aware exact \(0..2\) exit) \\
actual encoded suffix copy and tail frame & \statusproved \\
actual encoder--copy--decoder branch control & \statusproved\ (control only) \\
actual encoder produces the complete pure trace &
\statusproved\ (success-conditioned partial correctness) \\
actual decoder consumes the exact trace &
\statusproved\ (1024-symbol recovery, exact consumption, tail frame) \\
actual encoder--copy--decoder core inverse &
\statusproved\ (success-conditioned partial correctness) \\
full HBZ encode/decode &
\statusproved\ (production-wrapper success-conditioned partial correctness) \\
actual encoder success witness &
\statusproved\ (fixed all-6 Hoare failure exclusion; no termination/losslessness) \\
\bottomrule
\end{tabularx}
\end{center}
